\section{Compiling CFR onto Accelerators}
\label{sec:method}

The design follows one observation: once the game is fixed, the entire structure of
a CFR iteration is known at build time, and only the numbers change. We therefore split the solver into a one-time \emph{compilation} of the game into
flat arrays (\cref{sec:compiled}) and a per-iteration \emph{dataflow} over those
arrays whose operation sequence never changes (\cref{sec:dataflow}). The fixed
sequence makes the iteration a legal target for CUDA-graph capture
(\cref{sec:cudagraph}). \Cref{fig:archpipeline} shows the pipeline on an 11-node
example game with a chance root, two decisions per player, and six terminals, and
\cref{fig:archdataflow} reuses the same game. A layered verification design checks
the compiled schedule, graph replay, and evaluator (\cref{sec:correctness}).

% Arch-1: the game-to-dataflow compilation pipeline, drawn as objects rather than
% labeled boxes. Single-column, drawn at native size and fitted to the column width.
% Top: an 11-node example tree (1 chance root, 2 P1 decisions, 2 P2 decisions,
% 6 terminals) over alternating depth bands, with one-word glyph labels at the
% first occurrence of each node kind. Middle: thick mapping lines funnel every
% node into its breadth-first cell of the flat node array, above the values
% template. Bottom: the run-time submission timelines (eager vs captured graph).
\begin{figure}[!htb]
\begin{minipage}[b]{0.50\linewidth}\centering
\resizebox{\linewidth}{!}{%
\begin{tikzpicture}[
    font=\sffamily\small,
    dec1/.style={circle, draw=teal!60!black, fill=teal!45, minimum size=3.4mm, inner sep=0pt},
    dec2/.style={circle, draw=violet!60!black, fill=violet!40, minimum size=3.4mm, inner sep=0pt},
    chn/.style={diamond, draw=black!65, fill=black!22, minimum size=4.2mm, inner sep=0pt},
    ter/.style={rectangle, draw=black!70, fill=black!55, minimum size=2.5mm, inner sep=0pt},
    edge/.style={black!55, line width=0.9pt},
    flow/.style={-{Stealth[length=2.2mm]}, line width=1.0pt, black!60},
    fold/.style={-{Stealth[length=1.8mm]}, densely dashed, line width=0.8pt, black!55},
    mapline/.style={line width=1.0pt, opacity=0.80, shorten <=1.4pt},
    tick/.style={orange!80!black, line width=1.0pt},
    gtick/.style={orange!95!black, line width=1.6pt},
    note/.style={font=\sffamily\small, text=black!55},
    lanelabel/.style={font=\sffamily\small\bfseries, text=black!60}
]
% ================= build lane: tree over its flattened arrays =================
% alternating depth bands behind the tree
\foreach \y in {0, -1.7}
  \fill[blue!6] (0.10,\y+0.33) rectangle (7.95,\y-0.33);
% --- example game tree; left margin below stays clear for the compile label ---
\node[chn]  (c0)  at (4.60, 0.00) {};
\node[dec1] (p1a) at (2.60,-0.85) {};
\node[dec1] (p1b) at (6.60,-0.85) {};
\node[dec2] (q2a) at (1.70,-1.70) {};
\node[ter]  (t2a) at (3.50,-1.70) {};
\node[dec2] (q2b) at (5.70,-1.70) {};
\node[ter]  (t2b) at (7.60,-1.70) {};
\node[ter]  (t3a) at (1.30,-2.55) {};
\node[ter]  (t3b) at (2.30,-2.55) {};
\node[ter]  (t3c) at (5.20,-2.55) {};
\node[ter]  (t3d) at (6.20,-2.55) {};
\draw[edge, densely dashed] (c0) -- (p1a);
\draw[edge, densely dashed] (c0) -- (p1b);
\draw[edge] (p1a) -- (q2a); \draw[edge] (p1a) -- (t2a);
\draw[edge] (p1b) -- (q2b); \draw[edge] (p1b) -- (t2b);
\draw[edge] (q2a) -- (t3a); \draw[edge] (q2a) -- (t3b);
\draw[edge] (q2b) -- (t3c); \draw[edge] (q2b) -- (t3d);
% a compact standalone glyph legend keeps labels off the tree
\node[chn, minimum size=3.2mm] (legchance) at (0.55,0.62) {};
\node[note, anchor=west] at (0.98,0.62) {chance};
\node[dec1, minimum size=3.0mm] (legp1) at (2.30,0.62) {};
\node[note, anchor=west] at (2.73,0.62) {P1};
\node[dec2, minimum size=3.0mm] (legp2) at (3.48,0.62) {};
\node[note, anchor=west] at (3.91,0.62) {P2};
\node[ter, minimum size=2.4mm] (legterm) at (4.70,0.62) {};
\node[note, anchor=west] at (5.13,0.62) {terminal node};

% ================= flat arrays (BFS order, depth-segmented) =================
\def\cw{0.62}
\def\sx{0.72}
\def\slotb{-4.17}   % node strip bottom (top = -4.05)
\def\tmplb{-5.87}   % template strip bottom (top = -5.75)
% mapping fan: every tree node funnels into its BFS cell (compile = walk once).
% Each line first drops straight out of the tree, then converges inside the
% dedicated fan band so crossings never overlap the tree itself.
\foreach \n/\i/\c in {c0/0/black!50, p1a/1/teal!75, p1b/2/teal!75,
                      q2a/3/violet!75, t2a/4/black!50, q2b/5/violet!75,
                      t2b/6/black!50, t3a/7/black!50, t3b/8/black!50,
                      t3c/9/black!50, t3d/10/black!50}
  \draw[mapline, \c] (\n.south) -- (\n.south |- 0,-2.85)
    .. controls (\n.south |- 0,-3.25) and ({\sx+(\i+0.5)*\cw},-3.10) ..
    ({\sx+(\i+0.5)*\cw},\slotb+\cw);
\node[note, align=left, anchor=west] at (0.00,-3.02) {compile\\walk once};
% node strip: cells keep the color of their tree node, in BFS order
\node[note, anchor=east] at (\sx-0.10,\slotb+0.5*\cw) {nodes};
\foreach [count=\i from 0] \c in
  {black!22, teal!45, teal!45, violet!40, black!55, violet!40, black!55,
   black!55, black!55, black!55, black!55}{
  \draw[black!60, fill=\c] ({\sx+\i*\cw},\slotb) rectangle ++(\cw,\cw);
}
% chance slot carries no run-time work: overlay dashed hatch look
\draw[black!60, densely dashed, fill=white, fill opacity=0.55]
  (\sx,\slotb) rectangle ++(\cw,\cw);
% depth segment separators + labels
\foreach \s in {0, 1, 3, 7, 11}
  \draw[black!85, line width=1.2pt] ({\sx+\s*\cw},\slotb-0.05) -- ++(0,\cw+0.10);
\foreach \d/\a/\b in {0/0/1, 1/1/3, 2/3/7, 3/7/11}
  \node[note] at ({\sx+(\a+\b)*\cw/2},\slotb-0.26) {$d{=}\d$};
% per-depth index blocks: one gather/scatter block arrow per segment
\foreach \a/\b in {1/3, 3/7, 7/11}
  \draw[flow, line width=0.9pt, -{Stealth[length=1.8mm]}]
    ({\sx+\a*\cw+0.06},\slotb-0.56) -- ({\sx+\b*\cw-0.06},\slotb-0.56);
\node[note] at ({\sx+9*\cw},\slotb-0.86) {$D$ blocks};
% values template strip: only terminal cells are pre-filled
\node[note, anchor=east] at (\sx-0.10,\tmplb+0.5*\cw) {$\cfv_{\mathrm{tmpl}}$};
\foreach [count=\i from 0] \c in
  {white, white, white, white, black!55, white, black!55,
   black!55, black!55, black!55, black!55}{
  \draw[black!60, fill=\c] ({\sx+\i*\cw},\tmplb) rectangle ++(\cw,\cw);
}
\foreach \s in {0, 1, 3, 7, 11}
  \draw[black!85, line width=1.2pt] ({\sx+\s*\cw},\tmplb-0.05) -- ++(0,\cw+0.10);
% chance folding: the run-time-free chance slot melts into the template
\draw[fold] ({\sx+0.04},\slotb)
  .. controls (0.28,-4.60) and (0.28,-4.95) .. ({\sx+0.08},\tmplb+\cw)
  node[note, pos=0.5, left=0.6mm] {fold};

% ================= run lane: submission timelines =================
\begin{scope}[yshift=-0.10cm]
\def\wzero{1.66}
\def\itw{1.62}
\def\itg{0.14}
% three iteration windows, shaded alternately, shared by both rows
\foreach \i/\f in {0/black!5, 1/black!0, 2/black!5}
  \fill[\f] ({\wzero+\i*(\itw+\itg)},-7.30) rectangle ++(\itw,-1.15);
\node[note] at ({\wzero+0.5*\itw},-7.46) {one iteration};
% eager row: one launch tick per tensor op, dense within every iteration
\node[note, anchor=east] at (0.72,-7.78) {eager};
\draw[black!35] (\wzero,-7.78) -- ({\wzero+3*\itw+2*\itg},-7.78);
\foreach \i in {0,1,2}{
  \foreach \k in {0,...,8}{
    \draw[tick] ({\wzero+0.18+\i*(\itw+\itg)+\k*0.17},-7.65) -- ++(0,-0.26);
  }
}
% graph row: record once, then a single submission per iteration
\node[note, anchor=east] at (0.72,-8.25) {graph};
\draw[black!35] (\wzero,-8.25) -- ({\wzero+3*\itw+2*\itg},-8.25);
\node[draw=orange!80!black, fill=orange!25, rounded corners=1.5pt,
      minimum height=3.4mm, inner sep=1.6pt, anchor=east]
  at ({\wzero-0.08},-8.25) {record};
\foreach \i in {0,1,2}{
  \draw[gtick] ({\wzero+0.18+\i*(\itw+\itg)},-8.06) -- ++(0,-0.38);
}
\draw[flow, line width=0.9pt]
  ({\wzero+3*\itw+2*\itg},-8.25) -- ++(0.26,0);
\node[note, align=left, anchor=west] at (7.10,-8.25) {replay\\$\times T$};
% time axis hint
\draw[-{Stealth[length=1.6mm]}, black!40] (5.50,-8.65) -- (6.70,-8.65)
  node[note, pos=0.5, above=0.2mm] {time};
\end{scope}

% ================= lane frames =================
\begin{scope}[on background layer]
\node[draw=blue!45!black!30, rounded corners=3pt, inner sep=1.2mm, fill=blue!3,
      fit={(-0.22,0.88) (8.14,-6.35)}] (buildlane) {};
\node[draw=orange!70!black!30, rounded corners=3pt, inner sep=1.2mm, fill=orange!4,
      fit={(-0.22,-7.15) (8.14,-9.00)}] (runlane) {};
\end{scope}
\node[lanelabel, anchor=south west] at (buildlane.north west)
  {Build: once per game};
\node[lanelabel, anchor=south west] at (runlane.north west)
  {Run: every iteration};
\end{tikzpicture}}
\caption{The compilation pipeline on an 11-node example game. A build-time walk
flattens the tree into the depth-segmented node array and values template
(colored lines trace each node to its cell); at run time, eager execution
submits $c_1 + c_2 D$ framework operations per iteration (each one or more kernel
launches), graph replay one graph launch.}
\Description{Two-lane diagram. The build-time lane shows a small game tree with
a diamond chance root, colored decision nodes, and square terminals over
alternating depth bands; thick colored lines funnel every node into its
breadth-first cell of a flat array strip divided into depth segments, above a
values-template strip whose terminal cells are pre-filled; a short dashed arrow
folds the chance slot into that template. The run-time lane shows two aligned
timelines over three shaded iteration windows: the eager row has a dense comb
of launch ticks in every window, while the graph row has one record block
followed by a single tick per window and an arrow labeled replay times T.}
\label{fig:archpipeline}
\end{minipage}\hfill%

% Arch-2: one compiled CFR+ iteration on the same 11-node example game as
% Figure 2 (8 decision-edge slots, depth segments 1|2|4|4). Single column: the
% four phases stack top to bottom in execution order inside one rounded frame
% (the captured CUDA graph), with a left return edge closing the cycle. Every
% solver buffer is a strip of cells and every phase is arrows between specific
% cells. Arcs are colored by the edge's acting owner (gray chance, teal P1,
% violet P2) and staggered in height within each group; the reach buffer is one
% strip of split cells (teal top half = player 1 lane, violet bottom half =
% player 2 lane) because one index set advances both lanes. Formulas live in
% the body text, Alg. 1, and Table 1 (tab:opsteps).
\begin{minipage}[b]{0.46\linewidth}\centering
\resizebox{\linewidth}{!}{%
\begin{tikzpicture}[
    font=\sffamily\small,
    flow/.style={-{Stealth[length=2.6mm]}, line width=1.4pt, black!55},
    arr/.style={-{Stealth[length=1.6mm]}, line width=0.9pt, black!60},
    carc/.style={-{Stealth[length=1.6mm]}, line width=1.0pt, black!50},
    tarc/.style={-{Stealth[length=1.6mm]}, line width=1.0pt, teal!70!black},
    varc/.style={-{Stealth[length=1.6mm]}, line width=1.0pt, violet!70!black},
    sarr/.style={-{Stealth[length=1.5mm]}, line width=0.9pt, orange!75!black},
    note/.style={font=\sffamily\small, text=black!55},
    phaselabel/.style={font=\sffamily\small\bfseries, text=black!65},
    badge/.style={circle, fill=orange!85!black, text=white,
                  font=\sffamily\small\bfseries, inner sep=1.0pt}
]
\def\cw{0.58}
\def\sx{0.82}
% strip helpers: slot strips have 8 cells (+1 sentinel), node strips 11 cells
% with depth separators after cells 1, 3, 7 and dark terminal cells at
% BFS positions 4, 6, 7, 8, 9, 10.
\newcommand{\slotstrip}[3]{% x, y-bottom, fill
  \foreach \i in {0,...,7}
    \draw[black!60, fill=#3] ({#1+\i*\cw},#2) rectangle ++(\cw,\cw);
}
\newcommand{\nodestrip}[3]{% x, y-bottom, base fill (terminals dark)
  \foreach [count=\i from 0] \c in
    {#3, #3, #3, #3, black!55, #3, black!55, black!55, black!55, black!55, black!55}
    \draw[black!60, fill=\c] ({#1+\i*\cw},#2) rectangle ++(\cw,\cw);
  \foreach \s in {0, 1, 3, 7, 11}
    \draw[black!85, line width=1.0pt] ({#1+\s*\cw},#2-0.03) -- ++(0,\cw+0.06);
}
% dual-lane reach buffer as one strip of split cells: one index set advances
% both players, so each cell holds a teal (player 1) and violet (player 2) half
\newcommand{\dualstrip}[2]{% x, y-bottom
  \foreach \i in {0,...,10}{
    \fill[violet!25] ({#1+\i*\cw},#2) rectangle ++(\cw,0.5*\cw);
    \fill[teal!28]   ({#1+\i*\cw},{#2+0.5*\cw}) rectangle ++(\cw,0.5*\cw);
    \draw[black!60]  ({#1+\i*\cw},#2) rectangle ++(\cw,\cw);
  }
  \foreach \s in {0, 1, 3, 7, 11}
    \draw[black!85, line width=1.0pt] ({#1+\s*\cw},#2-0.03) -- ++(0,\cw+0.06);
}

% ================= phase 1 (top): regret matching =================
\node[badge] at (1.00,13.30) {1};
\node[phaselabel, anchor=west] at (1.22,13.30) {regret matching};
\node[note, anchor=east] at (\sx-0.10,12.74) {$\regret$};
\slotstrip{\sx}{12.45}{orange!30}
\node[note, anchor=east] at (\sx-0.10,11.64) {$\strat_{\mathrm{ext}}$};
\slotstrip{\sx}{11.35}{green!25}
% sentinel cell appended to the strategy strip
\draw[black!60, fill=black!45] ({\sx+8*\cw},11.35) rectangle ++(\cw,\cw);
\node[note] at ({\sx+8.5*\cw},11.13) {$\bot$};
% one arrow per slot, regret -> strategy; none enters the sentinel
\foreach \i in {0,...,7}
  \draw[arr] ({\sx+(\i+0.5)*\cw},12.40) -- ({\sx+(\i+0.5)*\cw},11.98);
\node[note, anchor=west] at ({\sx+8.3*\cw},12.19) {RM$^{+}$};
\draw[flow] (3.25,11.18) -- (3.25,10.78);

% ================= phase 2: forward reach =================
\node[badge] at (1.00,10.66) {2};
\node[phaselabel, anchor=west] at (1.22,10.66) {forward reach $\times D$};
\node[note, anchor=east] at (\sx-0.10,9.31) {$\reach_{1}$};
\node[note, anchor=east] at (\sx-0.10,9.00) {$\reach_{2}$};
\dualstrip{\sx}{8.85}
% parent -> child arcs above the strip, colored by the edge's acting owner;
% heights restart per owner group and grow within it. The small square at each
% apex is the gathered strategy factor: green = player slot, dark = sentinel.
\newcommand{\fwdarc}[6]{% parent, parent offset, child, height, arc style, factor
  \draw[#5] ({\sx+(#1+#2)*\cw},9.49)
    .. controls ({\sx+(#1+#2)*\cw},9.49+#4) and ({\sx+(#3+0.5)*\cw},9.49+#4) ..
    ({\sx+(#3+0.5)*\cw},9.49);
  \fill[#6, draw=black!70, line width=0.5pt]
    ({(\sx+(#1+#2)*\cw)/2+(\sx+(#3+0.5)*\cw)/2-0.08},{9.49+0.75*#4-0.08})
    rectangle ++(0.16,0.16);
}
% all ten edges are drawn (2 chance + 8 decision)
\fwdarc{0}{0.35}{1}{0.45}{carc}{black!45}
\fwdarc{0}{0.65}{2}{0.60}{carc}{black!45}
\fwdarc{1}{0.35}{3}{0.45}{tarc}{green!45}
\fwdarc{1}{0.65}{4}{0.60}{tarc}{green!45}
\fwdarc{2}{0.35}{5}{0.75}{tarc}{green!45}
\fwdarc{2}{0.65}{6}{0.90}{tarc}{green!45}
\fwdarc{3}{0.35}{7}{0.45}{varc}{green!45}
\fwdarc{3}{0.65}{8}{0.60}{varc}{green!45}
\fwdarc{5}{0.35}{9}{0.75}{varc}{green!45}
\fwdarc{5}{0.65}{10}{0.90}{varc}{green!45}
\draw[flow] (3.25,8.68) -- (3.25,8.28);

% ================= phase 3: backward values =================
\node[badge] at (1.00,8.05) {3};
\node[phaselabel, anchor=west] at (1.22,8.05) {backward values $\times D$};
\node[note, anchor=east] at (\sx-0.10,7.54) {$\cfv_{\mathrm{tmpl}}$};
\nodestrip{\sx}{7.25}{white}
\node[note, anchor=east] at (\sx-0.10,6.54) {$\cfv$};
\nodestrip{\sx}{6.25}{gray!15}
\draw[arr] ({\sx+4.5*\cw},7.20) -- node[note, right=0.4mm] {copy} ({\sx+4.5*\cw},6.88);
% child -> parent arcs under the value strip (reverse sweep), same edge colors
\newcommand{\bwdarc}[6]{% child, parent, parent offset, height, arc style, factor
  \draw[#5] ({\sx+(#1+0.5)*\cw},6.21)
    .. controls ({\sx+(#1+0.5)*\cw},6.21-#4) and ({\sx+(#2+#3)*\cw},6.21-#4) ..
    ({\sx+(#2+#3)*\cw},6.21);
  \fill[#6, draw=black!70, line width=0.5pt]
    ({(\sx+(#1+0.5)*\cw)/2+(\sx+(#2+#3)*\cw)/2-0.08},{6.21-0.75*#4-0.08})
    rectangle ++(0.16,0.16);
}
\bwdarc{1}{0}{0.35}{0.45}{carc}{black!45}
\bwdarc{2}{0}{0.65}{0.60}{carc}{black!45}
\bwdarc{3}{1}{0.35}{0.45}{tarc}{green!45}
\bwdarc{4}{1}{0.65}{0.60}{tarc}{green!45}
\bwdarc{5}{2}{0.35}{0.75}{tarc}{green!45}
\bwdarc{6}{2}{0.65}{0.90}{tarc}{green!45}
\bwdarc{7}{3}{0.35}{0.45}{varc}{green!45}
\bwdarc{8}{3}{0.65}{0.60}{varc}{green!45}
\bwdarc{9}{5}{0.35}{0.75}{varc}{green!45}
\bwdarc{10}{5}{0.65}{0.90}{varc}{green!45}
\draw[flow] (3.25,5.00) -- (3.25,4.68);

% ================= phase 4 (bottom): flat regret update =================
\node[badge] at (1.00,4.50) {4};
\node[phaselabel, anchor=west] at (1.22,4.50) {regret update, flat};
\node[note, anchor=east] at (\sx-0.10,3.99) {$\regret$};
\slotstrip{\sx}{3.70}{orange!30}
\node[note, anchor=east] at (\sx-0.10,2.59) {$\cfv$};
\nodestrip{\sx}{2.30}{gray!15}
% one scatter over all decision edges at once: child value cell -> its slot
\foreach \c/\s in {3/0, 4/1, 5/2, 6/3, 7/4, 8/5, 9/6, 10/7}
  \draw[sarr] ({\sx+(\c+0.5)*\cw},2.92)
    to[bend left=12] ({\sx+(\s+0.5)*\cw},3.66);

% ================= the iteration cycle =================
\draw[flow] (-0.30,4.30) -- (-0.30,12.50)
  node[note, midway, rotate=90, above=0.6mm] {next iteration};
% the whole cycle is one captured CUDA graph
\begin{scope}[on background layer]
\node[draw=black!30, rounded corners=4pt, fill=black!2, inner sep=1.2mm,
      fit={(-0.62,13.50) (7.38,1.92)}] (graphframe) {};
\end{scope}
\node[note, anchor=south east] at ($(graphframe.south east)+(-0.12,-0.02)$)
  {one CUDA graph};
\end{tikzpicture}}
\caption{One compiled CFR$^{+}$ iteration on the example game of
\cref{fig:archpipeline}, stacked in execution order inside the captured CUDA
graph (rounded frame). Arcs are colored by the acting owner of their edge (gray
chance, teal P1, violet P2); each apex square is the gathered strategy factor
(green) or the sentinel (dark).}
\Description{Four phases stacked top to bottom inside one rounded frame labeled
one CUDA graph. Phase one shows eight parallel arrows from an orange regret
strip to a green strategy strip with an extra dark sentinel cell that receives
no arrow. Phase two shows a reach strip of split teal and violet cells with
staggered arcs from parent cells to child cells, colored gray, teal, or violet
by the acting owner; each arc carries a small green square for a strategy
factor or a dark square for the sentinel. Phase three shows a template strip
copied into a value strip, with the same colored arcs folding child cells back
into parent cells. Phase four shows a gently curved fan of orange arrows
scattering child-value cells into the regret strip. A thick upward arrow on the
left labeled next iteration closes the cycle.}
\label{fig:archdataflow}
\end{minipage}
\end{figure}

\subsection{Compiled game representation}
\label{sec:compiled}

Compilation walks the game once, numbering nodes in breadth-first order so that
every parent precedes its children, and emits:
\begin{itemize}
    \item \textbf{Node arrays} of length $\numnodes$: owner (player 1/2, chance,
    terminal), depth, and a terminal payoff vector.
    \item \textbf{Edge arrays} of length $\numedges$: for every edge, its parent node,
    child node, and \emph{slot}, the index of the corresponding (information set,
    action) pair in a flat slot space of size $\numslots = \sum_{I} |\acts(I)|$.
    Chance edges carry their fixed probability and no slot.
    \item \textbf{Information-set arrays} of length $\numinfosets$ and $\numslots$:
    per-set action counts and slot offsets, so regret matching operates on flat
    vectors with segment reductions and never touches a ragged structure.
\end{itemize}
All solver state (accumulated regrets $\regret \in \R^{\numslots}$, strategy sums
$\bar{s} \in \R^{\numslots}$, and the iteration scratch below) lives in
preallocated device tensors. On solver construction, root reach entries are set to
one and every buffer is written before it is read. The scratch footprint is
$(5\numnodes + \numslots + \numinfosets + 1)$ floats plus the integer index arrays.
Measured peak GPU memory on the largest game is 183\,MiB allocated and 236\,MiB
reserved, including the CUDA-graph pool and evaluation scratch (\cref{tab:memory}).
That is about 570 bytes per tree node. Float scratch, persistent slot vectors, and
the int64 edge-index arrays all scale linearly in $\numnodes$, $\numedges$, and
$\numslots$, with $\numedges \approx \numnodes$ in trees. The same code path runs
unchanged on CPU and GPU.

\subsection{One iteration as a fixed dataflow}
\label{sec:dataflow}

A CFR iteration is regret matching, a forward reach pass, a backward value pass with
instantaneous-regret accumulation, and the update rule
(\crefrange{eq:reach}{eq:rm}). \Cref{fig:archdataflow} shows the four phases as they
execute. Three structural observations turn the two tree passes into a handful of
batched tensor operations each. Throughout, ``overwrite'' semantics are available
because the game is a tree. Every non-root node has exactly one incoming edge, so a
scatter along edges writes each destination exactly once per iteration, and buffers
never need zeroing within a solve.

\paragraph{(a) Static chance folding.}
Chance behavior is strategy-independent, so for every node $h$ the product of chance
probabilities on the root-to-$h$ path, $\pi_c(h)$, is a build-time constant. We fold
it into a \emph{values template}
\begin{equation}
    \cfv_{\mathrm{tmpl}}(h) \;=\;
    \begin{cases}
        \pi_c(h)\, u(h) & h \in \termnodes,\\
        0 & \text{otherwise},
    \end{cases}
    \label{eq:template}
\end{equation}
from which the backward pass starts every iteration. The invariant is that chance probabilities appear \emph{only} through
\cref{eq:template}. During the dynamic passes, every chance edge acts as multiplier
$1$ (via the sentinel slot below). The backward recurrence
$\cfv(h) = \sum_{(h,a,h')} \strat_{\mathrm{ext}}(h,a)\, \cfv(h')$
then returns at each node $h$ the \emph{chance-weighted} continuation value
$\sum_{z \succeq h} \pi_c(z)\, \reach^{\strat}_{1,2}(h \to z)\, u(z)$. There is
no double counting, because each terminal's full-path chance product enters exactly
once, at the template. This removes every chance multiplication from the
per-iteration forward pass. Chance edges remain in the depth blocks only as sentinel
reads.

\paragraph{(b) Depth-level execution blocks.}
The forward pass must respect topology, since a child's reach depends on its
parent's. In a tree, every child sits exactly one depth level below its parent, so
grouping edges by \emph{parent depth} yields a valid schedule at the coarsest
granularity a depth-synchronous schedule allows. All edges at one depth execute as a
single batched gather, multiply, and scatter, and the number of sequential steps
equals the tree depth. Across our suite this is 4--15 blocks per pass
(\cref{tab:suite}). A generic topological schedule over (stage $\times$ depth)
intersections produced up to 100 blocks on the same trees. The backward pass uses the same blocks
in reverse with \texttt{index\_add} scatters, which tolerate repeated parents within
a block.

\paragraph{(c) Sentinel slot and dual-lane reach buffer.}
Player $i$'s reach lane $\reach_i$ must multiply $i$'s own strategy entries and copy
through everything else:
\begin{equation}
    \reach_i(h') =
    \begin{cases}
        \reach_i(h)\, \strat(\mathrm{slot}(h,a)) & (h,a,h') \text{ owned by } i,\\
        \reach_i(h) & \text{otherwise}.
    \end{cases}
    \label{eq:lanes}
\end{equation}
We implement \cref{eq:lanes} branchlessly. The strategy vector is extended by one \emph{sentinel} slot pinned to $1.0$. For
each edge and each lane, a precomputed index points either at the edge's true slot
(own-player edges) or at the sentinel (all others). The reach buffer holds both lanes contiguously ($2\numnodes$ entries), and
per-block index arrays over the doubled range advance both lanes with a single
gather, multiply, and scatter. The same trick removes the last per-edge branch from regret accumulation. For a
player-$i$ edge $(h, a)$ with slot $q = \mathrm{slot}(h,a)$, the instantaneous
regret contribution is
\begin{equation}
    r(q) \;\mathrel{+}=\; s(e)\,\reach_{-i}(h)\,\big(\cfv(h') - \cfv(h)\big),
    \label{eq:instregret}
\end{equation}
where $e=(h,a,h')$, $s(e)=+1$ for player 1 and $-1$ for player 2,
$\cfv$ stores player-1 values, and $\reach_{-i}(h)$ is read from the
opponent's lane through a precomputed index array. Summing \cref{eq:instregret} over an information-set action with a slot-indexed
\texttt{index\_add} yields $\sum_{h\in I}s(e)\reach_{-i}(h)(\cfv(h')-\cfv(h))$,
the standard counterfactual regret. Chance reach is already in $\cfv$ through
\cref{eq:template} and is not multiplied again.

Per-block index arrays $P_2, C_2, S_2$ are the precomputed dual-lane parent/child/slot
arrays of block $b$, covering every edge (chance and other-player entries point at the
sentinel). $P_{\mathrm{val}}$, $C_{\mathrm{val}}$, $S_{\mathrm{val}}$ cover every edge of the
backward value pass; $P_{\mathrm{reg}}$, $C_{\mathrm{reg}}$, $S_{\mathrm{reg}}$,
$O_{\mathrm{reg}}$
cover player decision edges only, with $O_{\mathrm{reg}}$ indexing the opponent-reach lane at
the parent and $M_{\mathrm{reg}}$ the acting player's own lane there, used by the
averaging update. Because $\cfv$ stores player-1 values, a build-time sign
$s(e) \in \{\pm 1\}$ orients each decision edge's value difference toward its
acting player; $s(b)$ gathers these signs for the block's decision edges.
Chance edges and sentinel slots never receive regret updates.
$\strat_{\mathrm{ext}}$ denotes the strategy vector with the sentinel entry fixed at $1$.

\begin{algorithm}[t]
\caption{One compiled CFR$^{+}$ iteration under simultaneous updates.}
\label{alg:iteration}
\begin{algorithmic}[1]
\State $\strat \gets \mathrm{regret\_matching}(\regret)$ \Comment{flat over $\numslots$ slots; segment-normalized}
\For{$b = 1, \dots, D$} \Comment{forward, both lanes at once, \cref{eq:lanes}}
    \State $\reach[C_2(b)] \gets \reach[P_2(b)] \cdot \strat_{\mathrm{ext}}[S_2(b)]$
\EndFor
\State $\cfv \gets \cfv_{\mathrm{tmpl}}$ \Comment{chance-weighted payoffs, \cref{eq:template}}
\For{$b = D, \dots, 1$} \Comment{backward values}
    \State $\cfv.\mathrm{index\_add}\big(P_{\mathrm{val}}(b),\; \strat_{\mathrm{ext}}[S_{\mathrm{val}}(b)] \cdot \cfv[C_{\mathrm{val}}(b)]\big)$
    \State $\regret.\mathrm{index\_add}\big(S_{\mathrm{reg}}(b),\; s(b) \cdot \reach[O_{\mathrm{reg}}(b)] \cdot (\cfv[C_{\mathrm{reg}}(b)] - \cfv[P_{\mathrm{reg}}(b)])\big)$ \Comment{\cref{eq:instregret}, decision edges only}\label{alg:regretline}
\EndFor
\State $\bar{s} \mathrel{+}= w_t \cdot \reach_{\mathrm{own}} \cdot \strat$;\qquad
       $\regret \gets \max(\regret, 0)$ \Comment{update rule (here: CFR$^{+}$)}
\end{algorithmic}
\end{algorithm}

\cref{alg:iteration} summarizes the iteration, where $w_t$ is the averaging weight.
Alternating updates \cite{burch2019revisiting} run its body twice per iteration,
once per player, and restrict the regret update and the regret matching$^{+}$ clamp
to the acting player's slots. The player order is static, so every index array and
the captured graph stay the same. \Cref{tab:opsteps} lists the realized operation
schedule. One realization detail: the regret scatter of line~\ref{alg:regretline}
is hoisted out of the level loop and applied once, flat over all decision edges after
the value pass. The average-strategy accumulation shares the same flat scatter. The
parent and child values it reads are final by then, so the flat scatter accumulates
the same summands in a different grouping. Only the forward and value recurrences contribute
to the depth-proportional term. The result is a fixed sequence of eight Aten
\cite{paszke2019pytorch} operations per depth level plus a constant part: 64--152
operations per iteration across the suite. A reference implementation of the same
algorithm with per-(stage, depth) blocks, masked chance handling, and per-iteration
allocations needs 110--1{,}742 (\cref{tab:suite}). The Aten-operation count is
exactly reproducible and independent of machine load. It therefore serves both as a
regression gate (\cref{sec:correctness}) and as a benchmarking instrument.

% B3: the realized per-iteration operation schedule. Hand-authored from the
% solver's execution path; the two "per level" rows are the c2 D term of
% Proposition \ref{prop:depth}, everything else is the constant c1.
\begin{table}[!htb]
\centering
\caption{The realized operation schedule of one CFR$^{+}$ iteration
(\cref{fig:archdataflow}). $\mathrm{ia}$ abbreviates \texttt{index\_add},
$I(q)$ is the information set owning slot $q$, and the build-time sign
$s \in \{\pm 1\}$ orients each edge's regret toward its acting player. Only the
two per-level rows repeat with depth.}
\label{tab:opsteps}
{
\setlength{\tabcolsep}{3pt}
\begin{tabular}{@{}llr@{}}
\toprule
Realized operation & Data & Ops \\
\midrule
\multicolumn{3}{@{}l}{\emph{Regret matching (constant)}} \\
$\strat^{+} \gets \mathrm{clamp}(\regret, 0)$ & $[\numslots]$ & 1 \\
$\mathrm{tot} \gets 0$;\; $\mathrm{tot}.\mathrm{ia}(I(q), \strat^{+})$ & $[\numslots] {\to} [\numinfosets]$ & 2 \\
gather $\mathrm{tot}$ back to slots & $[\numinfosets] \to [\numslots]$ & 1 \\
$m \gets \mathrm{tot} > 0$;\; guard divisor via $\mathrm{where}$ & $[\numslots]$ & 2 \\
$\strat \gets \mathrm{where}(m, \strat^{+}/\mathrm{tot}, \mathrm{unif})$ & $[\numslots]$ & 2 \\
$\strat_{\mathrm{ext}}[{:}\numslots] \gets \strat$ \; (sentinel stays $1$) & $[\numslots{+}1]$ & 1 \\
\midrule
\multicolumn{3}{@{}l}{\emph{Forward reach, per level $\ell = 1, \dots, D$}} \\
$\reach[C_2(\ell)] \gets \reach[P_2(\ell)] \cdot \strat_{\mathrm{ext}}[S_2(\ell)]$ \; (both lanes) & $[2\numnodes]$ & $4 \cdot D$ \\
\midrule
\multicolumn{3}{@{}l}{\emph{Backward values, per level $\ell = D, \dots, 1$}} \\
$\cfv \gets \cfv_{\mathrm{tmpl}}$ \; (chance folded at build) & $[\numnodes]$ & 1 \\
$\cfv.\mathrm{ia}\big(P(\ell), \strat_{\mathrm{ext}}[S(\ell)] \cdot \cfv[C(\ell)]\big)$ & $[\numnodes]$ & $4 \cdot D$ \\
\midrule
\multicolumn{3}{@{}l}{\emph{Regret and average over decision edges (constant)}} \\
$\Delta \gets s \cdot (\cfv[C_{\mathrm{reg}}] - \cfv[P_{\mathrm{reg}}])$ & decision edges & 4 \\
$\regret.\mathrm{ia}\big(S_{\mathrm{reg}}, \reach[O_{\mathrm{reg}}] \cdot \Delta\big)$ & ${\to} [\numslots]$ & 3 \\
$\bar{s}.\mathrm{ia}\big(S_{\mathrm{reg}}, (w \cdot \reach[M_{\mathrm{reg}}] \cdot \strat_{\mathrm{ext}}[S_{\mathrm{reg}}])_{\mathrm{f64}}\big)$ & ${\to} [\numslots]$ & 6 \\
$\regret \gets \max(\regret, 0)$ \; (CFR$^{+}$) & $[\numslots]$ & 1 \\
$w \mathrel{{+}{=}} 1$ \; (inside the captured body) & scalar & 1 \\
\bottomrule
\end{tabular}}
\end{table}


\subsection{CUDA-graph execution}
\label{sec:cudagraph}

Even at 96 operations per iteration, a 434k-state game runs each GPU kernel for only
microseconds, so framework dispatch and kernel launches dominate wall-clock time. The
compiled dataflow executes eagerly at \eagerMsRange\,ms per iteration on an A100
across our suite, largely independent of tree size, the signature of a launch-bound
workload. Because the operation sequence, tensor shapes, and buffer
addresses are all fixed, the entire iteration qualifies for CUDA-graph capture
\cite{guide2020cuda}: record the kernel sequence once, then replay it with one graph
launch per iteration. CUDA graphs remove per-kernel launch and dispatcher overhead
and leave the kernels themselves unchanged, so the same kernels run back to back.
\Cref{app:profile} measures hand-fused kernels and \cref{app:compiler} lists the
compiler passes.

Three details make capture faithful to eager execution. First, CFR$^{+}$'s
linear-averaging weight $w_t = t$ changes every iteration, which a static graph
cannot see. We keep $w$ in a 0-dimensional device tensor that the captured body
increments in place after use. Before each replay batch the host refills it with the
exact starting index, so replay $k$ uses precisely $w = t_0 + k$ and matches eager
execution bit for bit. The counter is float32 in every reported run and represents
$w$ exactly up to $2^{24}$ iterations. Second, capture must be preceded by warmup
iterations on a side stream. In GPU-CFR these are genuine eager iterations that count
toward the iteration budget, and the accounting is preserved across arbitrarily split
solver calls. Third, replaying a graph whose buffers have been reallocated would
silently read stale memory. The solver stores the data pointers of all eight
persistent buffers at capture time and refuses to replay if any has changed. If
capture itself fails, for example on an unsupported driver, the solver warns once and
falls back to eager execution. The fallback runs the same kernels in the same order
and, within one process, computes the identical floats.

Graph replay cuts steady-state iteration time by \eagerGraphRange$\times$
depending on tree size (\cref{tab:matrix}, \cref{app:extra}). The two largest games
run at \graphMsTwoLargest\,ms per iteration. One-time capture cost
(0.037--0.301\,s) lands in the first solver call, and \cref{sec:results} reports
steady-state and end-to-end numbers separately.

\subsection{Layered verification}
\label{sec:correctness}

Verification proceeds at three levels. First, graph replay and compiled eager
execution run the same kernels in the same order and agree bitwise over 30 iterations
within one process.
Second, the compiled dataflow is compared with the pre-optimization implementation.
The two agree bitwise in float32 after 22 iterations on all eight benchmark games
(\cref{tab:suite}; CPU, fixed reduction order) and to $10^{-12}$ in float64 over 30
iterations. Third, independently written implementations provide tolerance-level
checks on values and exploitability.

GPU scatter reductions can accumulate in different orders across independent
processes. Near an indifferent action a last-bit difference can separate subsequent
trajectories: on the HUNL river subgame, 4 of 20 identical runs reach a second
exploitability cluster $0.6\%$ away. \Cref{app:correctness} traces this to
reduction order and quantifies its effect. The verification layers are:
\begin{enumerate}
    \item \textbf{Reference-oracle parity.} A complete copy of the solver from
    before optimization is preserved inside the test suite and compared as above,
    across $\{$CFR, CFR$^{+}\} \times \{$uniform, linear$\}$ averaging.
    \item \textbf{Independent reference solver.} A per-node Python implementation,
    sharing no code with the compiled path, must agree to $10^{-9}$ on the poker
    subgames.
    \item \textbf{Best-response oracles and analytic anchors.} An exact recursive best-response oracle \cite{johanson2011accelerating}, checked
    by brute-force strategy enumeration on small instances, validates the evaluator.
    Analytic cases pin absolute values.
    \item \textbf{Operation-count regression gates.} The per-iteration count of
    Aten operations is asserted against the affine schedule size $c_1 + c_2 D$ of
    \cref{sec:theory}, and graph replay is compared with eager execution over 30
    iterations. These deterministic checks complement wall-clock measurements.
\end{enumerate}
The full test inventory (499 tests) appears in \cref{app:correctness}.
